\documentclass[aspectratio=169,10pt]{beamer}
\usetheme[numbering=fraction,progressbar=frametitle]{metropolis}
\usepackage{fontspec}
\setsansfont{Latin Modern Sans}
\usepackage{amsmath,amssymb,mathtools,graphicx,booktabs}
\usepackage{polyglossia}
\setdefaultlanguage{english}
\setbeamertemplate{navigation symbols}{}
\setbeamerfont{frametitle}{size=\large}
\setbeamerfont{normal text}{size=\normalsize}
\setlength{\parskip}{.45em}
\hypersetup{pdfauthor={},pdfsubject={Lecture slides},pdfkeywords={},pdftitle={Complex numbers and geometry}}
\providecommand{\N}{\mathbb{N}}
\providecommand{\Z}{\mathbb{Z}}
\providecommand{\D}{\mathbb{D}}
\providecommand{\Q}{\mathbb{Q}}
\providecommand{\R}{\mathbb{R}}
\providecommand{\C}{\mathbb{C}}
\providecommand{\vect}[1]{\overrightarrow{#1}}
\providecommand{\vi}{\vec{\imath}}
\providecommand{\vj}{\vec{\jmath}}
\providecommand{\vk}{\vec{k}}
\providecommand{\norm}[1]{\left\lVert #1\right\rVert}
\providecommand{\vc}[2]{\begin{pmatrix}#1\\#2\end{pmatrix}}
\providecommand{\deter}[2]{\det\left(#1,#2\right)}
\providecommand{\abs}[1]{\left\lvert #1\right\rvert}
\providecommand{\dd}{\,\mathrm{d}}
\providecommand{\e}{\mathrm{e}}
\providecommand{\Card}{\operatorname{Card}}
\providecommand{\E}{\operatorname{E}}
\providecommand{\V}{\operatorname{V}}
\providecommand{\degre}{^{\circ}}
\providecommand{\un}[1]{\,\mathrm{#1}}
\providecommand{\sci}[2]{#1\times10^{#2}}
\providecommand{\pgcd}{\operatorname{PGCD}}
\providecommand{\ppcm}{\operatorname{PPCM}}
\title{Complex numbers and geometry}
\subtitle{Lecture and worked examples}
\date{}
\author{}
\begin{document}
\begin{frame}[plain]
\vfill
{\usebeamerfont{title}\inserttitle\par}
\medskip
{\usebeamerfont{subtitle}\insertsubtitle\par}
\medskip
{\usebeamercolor[fg]{progress bar}\rule{\textwidth}{.5pt}}
\vfill
\end{frame}
\begin{frame}{Objectives}
Interpret the modulus and an argument of \(\dfrac{c-a}{b-a}\); prove that points are collinear, that lines are perpendicular, the nature of a triangle; determine sets of points; determine the \(n\)-th roots of unity and use them for regular polygons.
\end{frame}
\begin{frame}{Introductory activity}
\begin{enumerate}
\setcounter{enumi}{0}\item Factorise \(z^3-1\) by \(z-1\), then solve \(z^3=1\) in \(\C\). Let \(j\) be the solution with positive imaginary part.

\item Plot the three solutions in the complex plane. Compute the distances between them. Which figure do you obtain?

\item Write \(j\) in exponential form. Check that \(j^2=\overline j\) and that \(1+j+j^2=0\). Interpret the last equality with the centroid.

\item Solve \(z^4=1\) and plot the solutions. Same question for \(z^6=1\), using \(z^6-1=\left(z^3-1\right)\left(z^3+1\right)\).

\end{enumerate}
\end{frame}
\begin{frame}{Introductory activity (2)}
\begin{enumerate}
\setcounter{enumi}{4}\item State a conjecture about the solutions of \(z^n=1\): their number, their modulus, their arguments, the figure they form.

\end{enumerate}
\end{frame}
\begin{frame}{Angles and lengths}
Let \(A\), \(B\), \(C\) be three points with affixes \(a\), \(b\), \(c\), with \(A\neq B\) and \(A\neq C\). Then
\[\abs{\dfrac{c-a}{b-a}}=\dfrac{AC}{AB}\qquad\text{and}\qquad\arg\left(\dfrac{c-a}{b-a}\right)\equiv\left(\vect{AB},\vect{AC}\right)\pmod{2\pi}.\]
\end{frame}
\begin{frame}{Angles and lengths}
The modulus of a quotient is the quotient of the moduli: \(\abs{c-a}=AC\) and \(\abs{b-a}=AB\). For the argument, \(\arg(b-a)\equiv\left(\vect u,\vect{AB}\right)\) and \(\arg(c-a)\equiv\left(\vect u,\vect{AC}\right)\), so by Chasles' relation \(\arg\dfrac{c-a}{b-a}\equiv\left(\vect u,\vect{AC}\right)-\left(\vect u,\vect{AB}\right)\equiv\left(\vect{AB},\vect{AC}\right)\pmod{2\pi}\).
\end{frame}
\begin{frame}{Angles and lengths}
With the same notation, let \(q=\dfrac{c-a}{b-a}\).

\begin{enumerate}
\setcounter{enumi}{0}
\item \(A\), \(B\), \(C\) are collinear if and only if \(q\) is real.

\item The lines \((AB)\) and \((AC)\) are perpendicular if and only if \(q\) is purely imaginary.

\item The triangle \(ABC\) is right-angled and isosceles at \(A\) if and only if \(q=i\) or \(q=-i\); it is equilateral if and only if \(q=\e^{i\pi/3}\) or \(q=\e^{-i\pi/3}\).

\end{enumerate}
\end{frame}
\begin{frame}{Angles and lengths}
\(q\) is real if and only if its argument is \(0\) or \(\pi\) modulo \(2\pi\), that is if \(\vect{AB}\) and \(\vect{AC}\) are collinear. \(q\) is purely imaginary (nonzero) if and only if its argument is \(\pm\dfrac\pi2\). For 3, a triangle that is right-angled and isosceles at \(A\) corresponds to \(AC=AB\) and \(\left(\vect{AB},\vect{AC}\right)\equiv\pm\dfrac\pi2\); an equilateral triangle to \(AC=AB\) and \(\left(\vect{AB},\vect{AC}\right)\equiv\pm\dfrac\pi3\).
\end{frame}
\begin{frame}{Angles and lengths}
Let \(A\) and \(B\) have affixes \(a\) and \(b\). The direct square \(ABCD\) has vertices \(c=b+i(b-a)\) and \(d=a+i(b-a)\): we go from \(\vect{AB}\) to \(\vect{BC}\) by a direct quarter turn. With \(a=0\) and \(b=1\): \(c=1+i\) and \(d=i\).
\end{frame}
\begin{frame}{\(n\)-th roots of unity}
Let \(n\ge1\) be an integer. An \(n\)-th root of unity is any complex number \(z\) such that \(z^n=1\). Their set is written \(\mathbb U_n\).
\end{frame}
\begin{frame}{\(n\)-th roots of unity}
For every integer \(n\ge1\), \(\mathbb U_n\) has exactly \(n\) elements:
\[\mathbb U_n=\left\{\e^{2ik\pi/n}\;:\;k\in\{0,\ldots,n-1\}\right\}.\]
\end{frame}
\begin{frame}{\(n\)-th roots of unity}
Let \(z\in\mathbb U_n\). Then \(\abs z^n=\abs{z^n}=1\) with \(\abs z\ge0\), so \(\abs z=1\) and \(z=\e^{i\theta}\) for some real number \(\theta\). The equality \(\e^{in\theta}=1\) is equivalent to \(n\theta\equiv0\pmod{2\pi}\), that is \(\theta=\dfrac{2k\pi}n\) with \(k\) an integer. Conversely, \(\left(\e^{2ik\pi/n}\right)^n=\e^{2ik\pi}=1\).
\end{frame}
\begin{frame}{\(n\)-th roots of unity (2)}
Two integers \(k\) and \(k'\) give the same number if and only if \(\dfrac{2(k-k')\pi}{n}\) is a multiple of \(2\pi\), that is \(k\equiv k'\pmod n\). Every integer \(k\) has the same remainder as an integer of \(\{0,\ldots,n-1\}\), and these \(n\) integers have distinct remainders: \(\mathbb U_n\) has exactly \(n\) elements.
\end{frame}
\begin{frame}{\(n\)-th roots of unity}
For \(n\ge3\), the images of the elements of \(\mathbb U_n\) are the vertices of a regular polygon with \(n\) sides inscribed in the unit circle, one of whose vertices is the point with affix \(1\). With \(\omega=\e^{2i\pi/n}\), we have \(\mathbb U_n=\{1,\omega,\omega^2,\dots,\omega^{n-1}\}\). Cases to know: \(\mathbb U_2=\{1,-1\}\), \(\mathbb U_3=\{1,j,j^2\}\) with \(j=\e^{2i\pi/3}\), \(\mathbb U_4=\{1,i,-1,-i\}\).
\end{frame}
\begin{frame}{\(n\)-th roots of unity}
For every integer \(n\ge2\), the sum of the \(n\)-th roots of unity is zero.
\end{frame}
\begin{frame}{\(n\)-th roots of unity}
\begin{columns}[T,onlytextwidth]
\begin{column}{.57\textwidth}
With \(\omega=\e^{2i\pi/n}\neq1\) (since \(n\ge2\)): \(1+\omega+\dots+\omega^{n-1}=\dfrac{1-\omega^n}{1-\omega}=0\).
\end{column}
\begin{column}{.4\textwidth}
\centering
\includegraphics[width=\linewidth,height=4.5cm,keepaspectratio]{figures/expertes-05-1-en.pdf}
\end{column}
\end{columns}
\end{frame}
\end{document}
